\section{Game-Theoretic Adversarial Analysis}

We model Byzantine attackers as rational agents maximizing attack impact subject to
detection probability constraints.
This analysis characterizes the \emph{optimal adaptive adversary}: an attacker that
knows the spectral filter and optimizes its gradient injections accordingly.

\begin{proposition}[Nash Equilibrium Strategy, empirically validated]
\label{thm:nash}
For rational Byzantine attackers that maximize
$\mathbb{E}[\text{attack impact}] - \lambda \cdot \mathbb{P}[\text{detection}]$
where $\lambda > 0$ is the cost of detection, the optimal strategy is:
\begin{equation}
g_{\mathrm{byz}} = \argmax_{\|g\| \leq \sigma\sqrt{d}} \left\{ \|g - \bar{g}\|^2 - \lambda \cdot \mathbb{P}[\mathrm{KS}(g) > \tau]\right\}
\end{equation}
Under this optimal adaptive adversary: (1) Below phase transition ($\sigma^2\beta^2 < 0.20$), our
protocol detects 96.7\% of attacks with 2.3\% false positive rate (measured empirically
across 500 attack trials); (2) Near phase transition ($0.20 \leq \sigma^2\beta^2 < 0.25$),
detection remains effective at 88.4\% with adaptive threshold calibration; (3) At or beyond
phase transition ($\sigma^2\beta^2 \geq 0.25$), detection becomes information-theoretically
impossible (Theorem~\ref{thm:phase_impossibility}).
We verify the empirical optimality of this strategy via exhaustive grid search over attacker
parameters; a formal Nash equilibrium proof remains future work.
\end{proposition}

\subsection{Differential Privacy Integration}

We demonstrate that combining our approach with $\varepsilon$-differential privacy ($\varepsilon=8$) via noise injection disrupts adversarial coordination while preserving the honest MP structure.
Table~\ref{tab:dp_extension} shows detection rates under DP.
Note that $\sigma^2\beta^2 \geq 0.25$ remains information-theoretically impossible for any gradient-only test (Theorem~\ref{thm:phase_impossibility}); the DP noise effectively reduces the \emph{attacker's} exploitable signal, lowering the \emph{effective} $\sigma^2\beta^2$ seen by the filter rather than pushing the fundamental limit.

\begin{table}[htbp]
\centering
\footnotesize
\caption{Differential Privacy Extension}
\label{tab:dp_extension}
\setlength{\tabcolsep}{3pt}
\begin{tabular}{|c|c|c|}
\hline
\textbf{Regime} & \textbf{No DP} & \textbf{$\varepsilon$-DP ($\varepsilon=8$)} \\
\hline
$\sigma^2\beta^2 < 0.20$ & 96.7\% & 94.5\% \\
$0.20 \leq \sigma^2\beta^2 < 0.25$ & 88.4\% & 85.2\% \\
$\sigma^2\beta^2 \geq 0.25$ & Impossible$^{\dagger}$ & 82.5\%$^{\dagger}$ \\
\hline
\end{tabular}
\end{table}

${}^{\dagger}$At $\sigma^2\beta^2 \geq 0.25$, Theorem~\ref{thm:phase_impossibility} establishes that no gradient-information-only test can distinguish Byzantine from honest gradients. The 82.5\% DP figure reflects a regime where DP noise shrinks the attacker's effective exploitable variance below the threshold, not a violation of the impossibility result.
